\section{Step 5: TCP Connection + Plain HTTP}

With a parsed host, port, and path we can finally speak HTTP. This step stays
on cleartext \texttt{http://} only; TLS arrives next. The flow is the classic
Berkeley sockets path: resolve $\rightarrow$ socket $\rightarrow$ connect
$\rightarrow$ write request $\rightarrow$ read response.

\subsection{Why These APIs?}

\begin{itemize}[leftmargin=1.4em]
    \item \texttt{getaddrinfo} replaces obsolete \texttt{gethostbyname} and
          supports IPv4 and IPv6 transparently (\texttt{AF\_UNSPEC}).
    \item Trying every result in the addrinfo list handles multi-homed hosts.
    \item HTTP/1.1 with \texttt{Connection: close} keeps the client simple:
          we read until EOF instead of parsing \texttt{Content-Length} or
          chunked encoding.
\end{itemize}

\subsection{Building the Request}

A minimal GET looks like:

\begin{lstlisting}[style=snippet,language=C]
GET /path HTTP/1.1\r\n
Host: example.com\r\n
User-Agent: ...\r\n
Accept: */*\r\n
Connection: close\r\n
\r\n
\end{lstlisting}

Note the blank line (\texttt{\textbackslash r\textbackslash n} alone) that
terminates the headers. Without it, many servers wait forever for more input.

\subsection{New Code --- Resolve, Connect, Request, Read}

\begin{lstlisting}[style=snippet,language=C]
struct addrinfo hints = {0};
hints.ai_family   = AF_UNSPEC;
hints.ai_socktype = SOCK_STREAM;
char ps[6];
snprintf(ps, sizeof(ps), "%u", u.port);

struct addrinfo *res = NULL;
if (getaddrinfo(u.host, ps, &hints, &res) != 0) {
    perror("getaddrinfo");
    return EXIT_FAILURE;
}

int s = -1;
for (struct addrinfo *a = res; a; a = a->ai_next) {
    s = socket(a->ai_family, a->ai_socktype, a->ai_protocol);
    if (s < 0) continue;
    if (connect(s, a->ai_addr, a->ai_addrlen) == 0) break;
    close(s); s = -1;
}
freeaddrinfo(res);
if (s < 0) { perror("connect"); return EXIT_FAILURE; }

char req[4096];
snprintf(req, sizeof(req),
    "GET %s HTTP/1.1\r\n"
    "Host: %s\r\n"
    "User-Agent: %s\r\n"
    "Accept: */*\r\n"
    "Connection: close\r\n"
    "\r\n",
    u.path, u.host, user_agent);
write(s, req, strlen(req));

char buf[8192];
ssize_t n;
while ((n = read(s, buf, sizeof(buf))) > 0)
    fwrite(buf, 1, n, stdout);
close(s);
\end{lstlisting}

If \texttt{u.https} is set we refuse for now---HTTPS needs OpenSSL next.

\FullSourceStart{urlgetstep05}
\begin{lstlisting}[style=oldcode,language=C,firstnumber=1]
/* urlget.c -- Step 5: TCP + plain HTTP */
#include <stdio.h>
#include <stdlib.h>
#include <string.h>
#include <stdbool.h>
#include <stdint.h>
\end{lstlisting}
\begin{lstlisting}[style=newcode,language=C,firstnumber=last]
#include <unistd.h>
#include <sys/types.h>
#include <sys/socket.h>
#include <netdb.h>
\end{lstlisting}
\begin{lstlisting}[style=oldcode,language=C,firstnumber=last]

typedef struct {
    bool https;
    char host[256];
    uint16_t port;
    char path[512];
} url_t;

static inline int parse_url(const char *arg, url_t *u)
{
    char buf[1024];
    if (strlen(arg) >= sizeof(buf)) return -1;
    strcpy(buf, arg);

    char *p = buf;
    u->https = false;

    if (strncmp(p, "https://", 8) == 0) {
        u->https = true; p += 8;
    } else if (strncmp(p, "http://", 7) == 0) {
        p += 7;
    }

    char *sep = strpbrk(p, "/?");
    char *hostport = p;

    if (sep) {
        if (*sep == '?') {
            size_t ql = strlen(sep);
            if (ql + 2 > sizeof(u->path)) return -1;
            u->path[0] = '/';
            memcpy(u->path + 1, sep, ql + 1);
        } else {
            if (strlen(sep) >= sizeof(u->path)) return -1;
            strcpy(u->path, sep);
        }
        *sep = '\0';
    } else {
        strcpy(u->path, "/");
    }

    char *host = u->host;
    size_t hmax = sizeof(u->host) - 1;
    uint16_t def = u->https ? 443 : 80;
    u->port = def;

    if (*hostport == '[') {
        char *cl = strchr(hostport, ']');
        if (!cl) return -1;
        size_t hl = cl - (hostport + 1);
        if (hl >= hmax) return -1;
        memcpy(host, hostport + 1, hl);
        host[hl] = '\0';
        char *aft = cl + 1;
        if (*aft == ':')
            u->port = (uint16_t)strtol(aft + 1, NULL, 10);
    } else {
        char *col = strrchr(hostport, ':');
        if (col) {
            size_t hl = col - hostport;
            if (hl >= hmax) return -1;
            memcpy(host, hostport, hl);
            host[hl] = '\0';
            u->port = (uint16_t)strtol(col + 1, NULL, 10);
        } else {
            if (strlen(hostport) >= hmax) return -1;
            strcpy(host, hostport);
        }
    }

    if (u->port == 0) u->port = def;
    return 0;
}

static inline void usage(const char *progname)
{
    fprintf(stderr,
        "Usage: %s [options] <URL>\n\n"
        "Options:\n"
        "  -A, --user-agent STR     Set User-Agent\n"
        "  -T, --tor [proxy]        Use Tor (default 127.0.0.1:9050)\n"
        "  -H, --header STR         Add custom header (repeatable)\n"
        "  -d, --data DATA          POST data (@file supported)\n"
        "  -o, --output FILE        Write output to file\n"
        "  -u, --user USER:PASS     Basic authentication\n"
        "  -I, --head               HEAD request\n"
        "      --no-location        Do not follow redirects\n"
        "      --max-redirects N    Max redirects (default: 10)\n",
        progname);
    exit(EXIT_FAILURE);
}

int main(int argc, char *argv[])
{
    const char *user_agent =
        "Mozilla/5.0 (X11; Linux x86_64; rv:128.0) "
        "Gecko/20100101 Firefox/128.0";
    const char *default_tor = "127.0.0.1:9050";

    const char *url_arg = NULL;
    const char *tor_proxy = NULL;
    bool use_tor = false;
    bool follow_redirects = true;
    int max_redirects = 10;
    bool head_request = false;
    char extra_headers[4096] = {0};
    char post_data[8192] = {0};
    char output_file[256] = {0};
    char auth_header[256] = {0};

    for (int i = 1; i < argc; i++) {
        if (strcmp(argv[i], "-A") == 0 ||
            strcmp(argv[i], "--user-agent") == 0) {
            if (i + 1 < argc) user_agent = argv[++i];
            else usage(argv[0]);
        }
        else if (strcmp(argv[i], "-T") == 0 ||
                 strcmp(argv[i], "--tor") == 0) {
            use_tor = true;
            if (i + 1 < argc && argv[i + 1][0] != '-' &&
                strchr(argv[i + 1], ':')) {
                tor_proxy = argv[i + 1];
                i++;
            }
        }
        else if (strcmp(argv[i], "-H") == 0 ||
                 strcmp(argv[i], "--header") == 0) {
            if (i + 1 < argc) {
                strncat(extra_headers, argv[++i],
                        sizeof(extra_headers) -
                        strlen(extra_headers) - 3);
                strcat(extra_headers, "\r\n");
            } else usage(argv[0]);
        }
        else if (strcmp(argv[i], "-d") == 0 ||
                 strcmp(argv[i], "--data") == 0) {
            if (i + 1 < argc) {
                const char *val = argv[++i];
                if (val[0] == '@') {
                    FILE *f = fopen(val + 1, "r");
                    if (f) {
                        fread(post_data, 1,
                              sizeof(post_data) - 1, f);
                        fclose(f);
                    }
                } else {
                    strncpy(post_data, val,
                            sizeof(post_data) - 1);
                }
            } else usage(argv[0]);
        }
        else if (strcmp(argv[i], "-o") == 0 ||
                 strcmp(argv[i], "--output") == 0) {
            if (i + 1 < argc)
                strncpy(output_file, argv[++i],
                        sizeof(output_file) - 1);
            else usage(argv[0]);
        }
        else if (strcmp(argv[i], "-u") == 0 ||
                 strcmp(argv[i], "--user") == 0) {
            if (i + 1 < argc)
                snprintf(auth_header, sizeof(auth_header),
                    "Authorization: Basic %s\r\n", argv[++i]);
            else usage(argv[0]);
        }
        else if (strcmp(argv[i], "-I") == 0 ||
                 strcmp(argv[i], "--head") == 0) {
            head_request = true;
        }
        else if (strcmp(argv[i], "--no-location") == 0) {
            follow_redirects = false;
        }
        else if (strcmp(argv[i], "--max-redirects") == 0 &&
                 i + 1 < argc) {
            max_redirects = atoi(argv[++i]);
        }
        else if (url_arg == NULL) {
            url_arg = argv[i];
        }
    }

    if (url_arg == NULL)
        usage(argv[0]);
    if (use_tor && tor_proxy == NULL)
        tor_proxy = default_tor;

    url_t u = {0};
    if (parse_url(url_arg, &u) != 0) {
        fprintf(stderr, "Failed to parse URL\n");
        return EXIT_FAILURE;
    }
\end{lstlisting}
\begin{lstlisting}[style=newcode,language=C,firstnumber=last]
    if (u.https) {
        fprintf(stderr, "HTTPS not implemented yet\n");
        return EXIT_FAILURE;
    }

    struct addrinfo hints = {0};
    hints.ai_family   = AF_UNSPEC;
    hints.ai_socktype = SOCK_STREAM;
    char ps[6];
    snprintf(ps, sizeof(ps), "%u", u.port);

    struct addrinfo *res = NULL;
    if (getaddrinfo(u.host, ps, &hints, &res) != 0) {
        perror("getaddrinfo");
        return EXIT_FAILURE;
    }

    int s = -1;
    for (struct addrinfo *a = res; a; a = a->ai_next) {
        s = socket(a->ai_family, a->ai_socktype, a->ai_protocol);
        if (s < 0) continue;
        if (connect(s, a->ai_addr, a->ai_addrlen) == 0) break;
        close(s); s = -1;
    }
    freeaddrinfo(res);
    if (s < 0) { perror("connect"); return EXIT_FAILURE; }

    char req[4096];
    snprintf(req, sizeof(req),
        "GET %s HTTP/1.1\r\n"
        "Host: %s\r\n"
        "User-Agent: %s\r\n"
        "Accept: */*\r\n"
        "Connection: close\r\n"
        "\r\n",
        u.path, u.host, user_agent);
    write(s, req, strlen(req));

    char buf[8192];
    ssize_t n;
    while ((n = read(s, buf, sizeof(buf))) > 0)
        fwrite(buf, 1, n, stdout);
    close(s);

    (void)use_tor; (void)tor_proxy;
    (void)follow_redirects; (void)max_redirects;
    (void)head_request; (void)extra_headers;
    (void)post_data; (void)output_file; (void)auth_header;
\end{lstlisting}
\begin{lstlisting}[style=oldcode,language=C,firstnumber=last]
    return 0;
}
\end{lstlisting}
\FullSourceStop

Test with a plain HTTP site (many sites redirect to HTTPS---that is fine; you
will still see the redirect response headers and body):

\begin{lstlisting}[style=snippet,language=bash]
./urlget http://example.com/
\end{lstlisting}
